\section{From ESS Requirements to Chain Steps}
\label{sec:ess}

\subsection{The ESS-Autocorrelation Relationship}

G.4~\citep{g4paper} provides empirical estimates of lag-1 Hamming
autocorrelation $\rho_1$ for \recom{} chains across six states.
The ESS for a chain of $n$ steps with geometric autocorrelation decay
$\rho_k \approx \rho_1^k$ is:
\[
  n^* = \frac{n}{1 + 2\sum_{k=1}^\infty \rho_1^k}
       = \frac{n(1 - \rho_1)}{1 + \rho_1}.
\]

Rearranging: to achieve $n^*$ effective samples, the required chain length is:
\[
  n = n^* \cdot \frac{1 + \rho_1}{1 - \rho_1}.
\]

\subsection{State-Specific Calculation}

From G.4 Table~4~\citep{g4paper}, the empirical $\rho_1$ values are:

\begin{table}[h]
\centering
\caption{Lag-1 Hamming autocorrelation $\rho_1$ from G.4, and chain steps
         required to achieve $n^* = 385$ effective samples for 80\% power
         at $\alpha = 0.05$ to detect a 10pp partisan advantage.}
\label{tab:autocorr}
\begin{tabular}{lcccc}
\toprule
State & $k$ & $\rho_1$ & $(1+\rho_1)/(1-\rho_1)$ & Steps for $n^* = 385$ \\
\midrule
NC & 14 & 0.87 & 14.4 & $5{,}544$ \\
WI &  8 & 0.85 & 12.3 & $4{,}736$ \\
GA & 14 & 0.88 & 15.7 & $6{,}044$ \\
PA & 17 & 0.89 & 17.2 & $6{,}622$ \\
TX & 38 & 0.92 & 24.0 & $9{,}240$ \\
CA & 52 & 0.94 & 32.7 & $12{,}589$ \\
\bottomrule
\end{tabular}
\end{table}

For the six G.4 study states, the required chain lengths for 80\% power
range from approximately 5,000 steps (WI) to 13,000 steps (CA).
This is the ESS-based requirement; the G.4 convergence requirement
(steps to $\rhat < 1.1$) is an additional constraint.

\subsection{The Binding Constraint: Convergence vs.\ Power}

Two constraints must both be satisfied:
\begin{itemize}
\item \textbf{Convergence}: Steps $\geq t_{\rm mix}$ (from G.4 Table~6),
  where $t_{\rm mix}$ is the steps to $\rhat < 1.1$.
\item \textbf{Power}: Steps $\geq n^* \cdot (1+\rho_1)/(1-\rho_1)$.
\end{itemize}

For smaller states (NC, WI), the power constraint is binding:
NC requires 5,544 steps for power but only 5,000 for convergence.
For larger states (TX, CA), the convergence constraint is binding:
CA requires 20,000 steps for convergence but only 12,589 for power.

The minimum chain length is $\max(t_{\rm mix}, n_{\rm power})$.

\subsection{Extrapolation to All 50 States}

For states not in G.4's study set, we interpolate $\rho_1$ from the
observed relationship between $k$ (delegation size) and $\rho_1$:
\[
  \rho_1(k) \approx 0.83 + 0.0025 \cdot k \quad \text{(empirical fit to G.4 data)}.
\]
This gives $\rho_1(1) = 0.832$ for single-district states (where no
redistricting is needed) and $\rho_1(52) = 0.96$ for California.

The convergence time $t_{\rm mix}$ scales approximately as:
\[
  t_{\rm mix}(k) \approx 200 \cdot k \quad \text{(empirical, from G.4 Table 6)}.
\]
This gives $t_{\rm mix}(14) \approx 2{,}800$ (consistent with G.4's NC estimate
of $\approx 2{,}000$), $t_{\rm mix}(38) \approx 7{,}600$ (consistent with G.4's
TX estimate of $\approx 10{,}000$), and $t_{\rm mix}(52) \approx 10{,}400$
(consistent with G.4's CA estimate of $\approx 20{,}000$ with the higher
$\rho_1$).

The minimum step count for each state is:
\[
  n_{\rm min}(k) = \max\!\left(
    200k,\;
    385 \cdot \frac{1 + \rho_1(k)}{1 - \rho_1(k)}
  \right).
\]
